\section{Приказ апликације}

У овом поглављу апликација је представљена кроз приказ њених екрана, редоследом којим се корисник са њима сусреће при типичној употреби. За сваки екран дат је кратак опис садржаја и могућих интеракција, праћен снимком екрана. Кориснички интерфејс изведен је у складу са Материјал дизајн (енгл. \textit{Material Design}) смерницама, уз употребу компоненти библиотеке \texttt{Material Components} за Андроид.

\subsection{Почетни екран}

Почетни екран представља улазну тачку апликације. У горњем делу приказани су логотип у заобљеном блоку, назив \texttt{AR Pametne Učionice} и кратак опис који кориснику објашњава да се усмеравањем камере на маркер учионице изнад њега приказују живи подаци о просторији. Испод описа налази се текстуално поље за адресу сервера, унапред попуњено вредношћу \texttt{http://192.168.0.49:8080}; свака измена се одмах трајно чува у дељеним подешавањима, па се при следећем покретању не мора уносити поново. На дну екрана су два дугмета: главно дугме \texttt{Pokreni AR pregled} отвара АР екран, док дугме \texttt{Uputstvo} води на екран са упутством.

\snimakekrana{slike/02-pocetni-ekran.png}{Почетни екран апликације}

\subsection{Екран са упутством}

Екран са упутством садржи четири нумерисана корака приказана у виду засебних картица. Први корак упућује корисника да пронађе маркер на вратима учионице, други да камеру телефона усмери ка маркеру, трећи објашњава да се изнад маркера приказује панел са живим подацима о просторији, а четврти да боје индикатора показују статус измерених вредности --- зелена да је све у реду, жута упозорење, а црвена критично стање.

\snimakekrana{slike/03-uputstvo.png}{Екран са упутством за коришћење}

\subsection{АР екран --- препознавање маркера}

Након покретања АР прегледа отвара се екран на коме слика са камере испуњава целу површину приказа. На врху се налази тамна трака са дугметом за затварање и насловом апликације, а испод ње статусна порука у облику пилуле која у почетном стању гласи \texttt{Pronađi marker učionice.} и која корисника током рада води кроз процес. У овој фази корисник полако помера телефон док камера не обухвати одштампани маркер на вратима учионице.

\snimakekrana{slike/04-ar-prepoznavanje.png}{АР екран пре препознавања маркера}

\subsection{АР екран --- усидрени панел}

Чим ARCore препозна маркер, изнад његове горње ивице појављује се усидрени панел са подацима о учионици, а статусна порука се мења у \texttt{Učionica 101 — podaci uživo.}. Панел приказује назив просторије, ред заузетости са називом текућег часа и временом до којег је учионица заузета, три реда са температуром, буком и угљен-диоксидом уз индикаторе обојене према статусу и препоруку сервера. Панел остаје причвршћен за маркер док се корисник креће: може му се прићи, одмаћи се или га посматрати са стране, а садржај се освежава на сваке три секунде свежим подацима са сервера.

\snimakekrana{slike/05-ar-panel.png}{Усидрени панел изнад маркера са живим подацима}

\subsection{Демонстрација са контролном таблом}

Завршни део демонстрације повезује све три компоненте система. Презентер на контролној табли помери клизач --- на пример, подигне вредност угљен-диоксида изнад критичног прага --- и промена се после задршка од 300 милисекунди шаље серверу. Већ при следећем преузимању, најкасније три секунде касније, АР панел на телефону приказује нову вредност, индикатор квалитета ваздуха прелази у црвено и исписује се препорука \texttt{Hitno provetriti prostoriju.}, без иједне акције на самом телефону.

\snimakekrana{slike/06-kontrolna-tabla-demo.png}{Промена на контролној табли одражена на АР панелу}

Приказани екрани заједно чине заокружену целину: од почетног екрана и упутства, преко препознавања маркера, до живог панела који у реалном времену одражава стање паметног окружења. Демонстрациони ток --- клизач на контролној табли, упис на серверу, освежен панел у проширеној стварности у року од три секунде --- показује да је архитектура спремна и за стварну примену: довољно је да уместо контролне табле исте крајње тачке позивају прави сензори уграђени у учионице, без иједне измене у АР апликацији. Као могућа проширења у будућем раду издвајају се повезивање стварних сензора температуре, буке и угљен-диоксида преко постојећег интерфејса, додавање маркера за већи број учионица и зграда, као и интеракција на самом панелу, попут резервације слободне учионице додиром.
